% ==============================================================================
%  MAPLE: Manifold-Activated Pruning of Latent Experts for Sub-2-BPW MoE LLMs
%  Submission Target: 9th Workshop on Accelerated Machine Learning (AccML 2027)
%  Co-located with HiPEAC 2027, Glasgow, United Kingdom
%  Review Mode: Single-blind (Author names & affiliations explicitly visible)
%  Format: IEEE Double-Column Conference (4–8 pages)
% ==============================================================================

\documentclass[conference]{IEEEtran}

\usepackage[utf8]{inputenc}
\usepackage[T1]{fontenc}
\usepackage{amsmath,amssymb,amsfonts}
\usepackage{algorithmic}
\usepackage{graphicx}
\usepackage{textcomp}
\usepackage{xcolor}
\usepackage{booktabs}
\usepackage{url}
\usepackage{listings}
\usepackage{cite}
\usepackage{microtype}

\lstset{
  basicstyle=\ttfamily\scriptsize,
  keywordstyle=\color{blue}\bfseries,
  commentstyle=\color{gray},
  stringstyle=\color{teal},
  breaklines=true,
  frame=single,
  tabsize=2
}

\newcommand{\MAPLE}{\textsc{MAPLE}}
\newcommand{\bpw}{\text{BPW}}
\newcommand{\tps}{\text{TPS}}
\newcommand{\R}{\mathbb{R}}

\begin{document}

\title{\MAPLE: Manifold-Activated Pruning of Latent Experts for Sub-2-BPW Mixture-of-Experts Language Models}

\author{
  \IEEEauthorblockN{Aman Sachan}
  \IEEEauthorblockA{\textit{Independent Systems Research} \\
  New Delhi, India \\
  amansachan92905@gmail.com}
}

\maketitle

\begin{abstract}
We present \MAPLE{} (\textbf{M}anifold-\textbf{A}ctivated \textbf{P}runing of \textbf{L}atent \textbf{E}xperts), a post-training structural compression methodology for fine-grained Mixture-of-Experts (MoE) large language models that achieves sub-2-bit-per-weight effective parameter density while retaining high task-level semantic fidelity. While global multilingual pre-training necessitates massive expert populations (e.g., 128--256 experts), typical edge deployment domains draw upon a surprisingly compact manifold of routed expert capacity. We formalize this sparsity by profiling empirical forward-pass expert activations over a representative domain corpus. Each expert is evaluated via a composite metric combining activation frequency and input feature energy: $s_{e,l} = C_{e,l} \cdot \sqrt{\bar{\sigma}_{e,l} + \epsilon}$. Applied to Gemma~4~26B-A4B (128 routed experts, 30 layers, 9.92~GB in IQ2\_XXS quantization), \MAPLE{} excises 80 dormant experts per layer to assemble an optimized 48-expert container of 4.75~GB (effective 1.48 bits per parameter, 2.09$\times$ compression). The compressed model achieves 100\% GPU VRAM residency on a consumer 6~GB laptop GPU (RTX~4050), accelerating generation by 6.57$\times$ from 6.9 to 45.4 tokens/second by eliminating PCIe memory-paging bottlenecks. Standardized task evaluation across code synthesis, data structure implementation, scientific deduction, and formal logic demonstrates a 93.8\% mean keyword fidelity retention. We document three subtle failure modes overcome during design: proxy-embedding routing estimation, cascaded normalization suppression, and DRY sampling interference.
\end{abstract}

\begin{IEEEkeywords}
Mixture of Experts, Model Compression, Edge Inference, Structured Pruning, GGUF, Quantization, High-Throughput Systems.
\end{IEEEkeywords}

\section{Introduction}

Mixture-of-Experts (MoE) architectures~\cite{jiang2024mixtral,dai2024deepseekmoe,team2025gemma} have become the foundational paradigm for scaling modern large language models, decoupling total parameter capacity from per-token compute by routing inputs through a sparse subset of $k$ active experts out of $E$ total candidates ($k \ll E$). While this design maintains sub-linear FLOP counts during generation, running frontier MoE models on consumer edge hardware remains fundamentally bottlenecked by GPU Video RAM (VRAM) capacity.

Even when quantized to extreme low-bit representations such as 2-bit quantization (e.g., Gemma~4~26B-A4B at IQ2\_XXS), the resulting model footprint exceeds 9.92~GB. On consumer laptops and workstations equipped with common 6~GB or 8~GB GPUs (such as the NVIDIA RTX 4050 or RTX 3060), the model cannot fit into high-bandwidth VRAM. Consequently, inference runtimes must offload layers to system RAM over the PCIe bus. This offloading creates severe memory-bandwidth stalls, causing throughput to collapse from interactive conversational speeds to single digits (6.9 tokens/second on an RTX 4050 Laptop GPU). In edge computing, achieving 100\% in-VRAM execution represents a discrete operational cliff: crossing below the physical VRAM threshold transforms a model from an unusable system into a real-time local copilot.

Existing compression literature is largely bifurcated into:
\begin{enumerate}
\item \textit{Numerical Quantization:} Techniques such as GPTQ~\cite{frantar2022gptq}, AWQ~\cite{lin2024awq}, and SpQR~\cite{dettmers2023spqr} compress weights from 16-bit floats to 4-bit, 3-bit, or 2-bit integers. However, pushing numerical precision below 2 bits per weight yields steep quantization noise, gradient instability, and catastrophic perplexity explosion.
\item \textit{Structural Pruning:} Unstructured and semi-structured weight pruning methods (e.g., SparseGPT~\cite{frantar2023sparsegpt}) typically require specialized sparse matrix multiplication kernels that yield minimal hardware acceleration on consumer GPUs without custom tensor cores.
\end{enumerate}

Fine-grained MoE models provide an orthogonal, unexplored dimension for compression: \emph{expert capacity allocation}. During pre-training, an expansive expert pool ($E = 64\text{--}128$) is required to capture hundreds of human languages, diverse programming dialects, and encyclopedic facts. However, edge deployments are rarely omni-domain. In practice, a local deployment is dedicated to a specific operational subdomain, such as English technical reasoning, algorithmic development, and conversational assistance. Within this domain, a substantial fraction of experts remain completely unselected.

In this paper, we introduce \textbf{\MAPLE{}} (\textbf{M}anifold-\textbf{A}ctivated \textbf{P}runing of \textbf{L}atent \textbf{E}xperts), an empirical forward-pass profiling and structural pruning framework designed specifically for fine-grained MoE architectures in the GGUF ecosystem~\cite{llama_cpp}. \MAPLE{} captures the empirical activation manifold of the target domain, computes a joint frequency-energy saliency score for every expert across all layers, and directly excises dormant experts from quantized GGUF binaries without requiring retraining, gradient access, or dequantization.

Applied to Gemma~4~26B-A4B, \MAPLE{} compresses the model from 9.92~GB to 4.75~GB—achieving an effective density of **1.48 bits per parameter** (a 2.09$\times$ reduction over an already 2-bit quantized base). Crucially, the 4.75~GB model fits entirely within the 6141~MiB VRAM budget of an RTX 4050 Laptop GPU, accelerating generation by **6.57$\times$ (from 6.9 to 45.4 TPS)** while retaining **93.8\% semantic task fidelity**.

\section{The MAPLE Framework}

\subsection{MoE Architectural Mechanics}
In a standard MoE transformer block, the conventional feed-forward network (FFN) is replaced with a routing gate and $E$ parallel expert networks. Given an incoming hidden state $\mathbf{h} \in \R^d$, the router computes unnormalized logits $\mathbf{g} = \mathbf{h} W_\text{gate}$, where $W_\text{gate} \in \R^{d \times E}$. The top-$k$ experts with the highest routing scores are selected:
\begin{equation}
\mathcal{T}_k(\mathbf{h}) = \text{TopK}(\mathbf{g}, k)
\end{equation}
The gate probabilities are computed via a softmax over the selected top-$k$ logits:
\begin{equation}
p_e(\mathbf{h}) = \frac{\exp(g_e)}{\sum_{j \in \mathcal{T}_k(\mathbf{h})} \exp(g_j)}, \quad \forall e \in \mathcal{T}_k(\mathbf{h})
\end{equation}
The block output is the weighted sum of the active experts' outputs plus any shared dense projection:
\begin{equation}
\mathbf{y} = \text{FFN}_\text{dense}(\mathbf{h}) + \sum_{e \in \mathcal{T}_k(\mathbf{h})} p_e(\mathbf{h}) \cdot \text{FFN}_e(\mathbf{h})
\end{equation}
In Gemma~4~26B-A4B, $E = 128$ routed experts exist per block across $L = 30$ layers, with $k = 8$ experts active per token. The routed expert weights (\texttt{ffn\_down\_exps} and \texttt{ffn\_gate\_up\_exps}) constitute 8.24~GB of the total 9.92~GB container size.

\subsection{Domain Calibration Corpus}
To discover which experts form the operative manifold, we construct a calibration dataset $\mathcal{D}_\text{calib}$ covering the target deployment distribution. For our implementation, we assembled a 1,979-token multi-domain corpus structured under the model's native conversational template (\texttt{<|turn>user ... <turn|> <|turn>model}):
\begin{itemize}
\item \textbf{Algorithmic Code Synthesis (30\%):} Recursive functions, dynamic programming, sorting routines, binary search trees, and Python decorators.
\item \textbf{Mathematical Deduction (20\%):} Matrix algebra, combinatorial probability, and continuous optimization.
\item \textbf{Scientific Reasoning (20\%):} Thermodynamics, Rayleigh scattering, wave mechanics, and chemical kinetics.
\item \textbf{Formal Deductive Logic (15\%):} Syllogisms, transitive inference chains, and set inclusion.
\item \textbf{Structured Discourse (15\%):} Technical summarization and precise instruction adherence.
\end{itemize}

\subsection{Empirical Forward-Pass Profiling}
Rather than approximating expert importance using static weight statistics, \MAPLE{} captures the live forward-pass activation statistics using \texttt{llama-imatrix}. As tokens traverse the network, the profiler tracks two empirical tensors for each layer $l \in [0, L-1]$:

\begin{enumerate}
\item \textbf{Activation Frequency Vector ($C_{e,l} \in \mathbb{N}^E$):} The exact count of tokens routed to expert $e$:
\begin{equation}
C_{e,l} = \sum_{\tau \in \mathcal{D}_\text{calib}} \mathbb{I}(e \in \mathcal{T}_k(\mathbf{h}_{\tau,l}))
\end{equation}
\item \textbf{Input Energy Sum Matrix ($\Sigma_{e,l} \in \R^{E \times d_\text{ff}}$):} The accumulated squared magnitude of the input features across all $d_\text{ff} = 704$ intermediate projection channels:
\begin{equation}
\Sigma_{e,l,j} = \sum_{\tau \in \mathcal{D}_\text{calib}: e \in \mathcal{T}_k(\mathbf{h}_{\tau,l})} (h_{\tau,l,j})^2
\end{equation}
\end{enumerate}
From the energy sum, we compute the mean channel energy:
\begin{equation}
\bar{\sigma}_{e,l} = \frac{1}{d_\text{ff}} \sum_{j=1}^{d_\text{ff}} \Sigma_{e,l,j}
\end{equation}

\subsection{The MAPLE Composite Saliency Score}
A naive frequency count $C_{e,l}$ ignores the magnitude of the signal processed by the expert: an expert called frequently on low-norm padding tokens may contribute minimally to the residual stream. Conversely, sorting purely by raw energy $\bar{\sigma}_{e,l}$ is fragile against stochastic outliers.

\MAPLE{} defines a balanced composite scoring function:
\begin{equation}
s_{e,l} = C_{e,l} \cdot \sqrt{\bar{\sigma}_{e,l} + \epsilon}
\end{equation}
where $\epsilon = 10^{-8}$ prevents numerical instability. 

For each layer $l$, the $E=128$ experts are sorted by $s_{e,l}$ descending, and the top $k'=48$ expert indices are extracted into a layer selection map $\mathcal{S}_l \subset \{0, \ldots, 127\}$ with $|\mathcal{S}_l| = 48$.

\subsection{Empirical Activation Mass Coverage}
We define the layer-wise empirical activation coverage as:
\begin{equation}
\text{Coverage}(\mathcal{S}_l) = \frac{\sum_{e \in \mathcal{S}_l} C_{e,l}}{\sum_{e=0}^{E-1} C_{e,l}} \times 100\%
\end{equation}
Empirically, selecting $k'=48$ experts captures an average of **79.88\%** of the total forward-pass activation mass across all 30 layers (ranging from 71.2\% to 92.4\% per layer). The discarded 80 experts account for 62.5\% of the expert parameter volume but contribute only 20.12\% of forward activations on the deployment distribution.

\subsection{Zero-Dequantization Binary Slicing}
To assemble the target model without introducing quantization noise or requiring retraining, \MAPLE{} streams the source GGUF binary and performs direct axis-0 index slicing on the quantized tensors:
\begin{enumerate}
\item \texttt{blk.\{l\}.ffn\_gate\_inp.weight}: The router gate matrix $[E, d]$ is sliced along axis 0 to $[\mathcal{S}_l, :]$, yielding shape $[48, d]$.
\item \texttt{blk.\{l\}.ffn\_down\_exps.weight}: Sliced along expert axis 0 to retain $[\mathcal{S}_l, \ldots]$ in native IQ2\_XXS format.
\item \texttt{blk.\{l\}.ffn\_gate\_up\_exps.weight}: Sliced along expert axis 0 to retain $[\mathcal{S}_l, \ldots]$.
\item \texttt{blk.\{l\}.ffn\_down\_exps.scale}: Quantization scales are sliced to $[\mathcal{S}_l, \ldots]$ with scale multipliers locked at $1.0\times$.
\item \textbf{Invariant Tensors:} Attention query, key, value, and output matrices, input/output embeddings, normalization layers, and shared dense FFN blocks are copied directly without byte modification.
\end{enumerate}

\section{Analysis of Failure Modes}

Achieving high semantic fidelity at 1.48 BPW required resolving three distinct failure modes encountered during development. We document these mechanisms to prevent recurrent pitfalls in MoE pruning.

\subsection{Failure Mode 1: The Ghost Router (Proxy-Embedding)}
\textbf{Symptom:} Repetitive token generation loops (e.g., infinite sequences of \texttt{nums\_nums\_nums...}) and total inability to complete logical syllogisms.

\textbf{Root Cause:} An early attempt bypassed forward passes by estimating routing logits from the static token embedding matrix $W_\text{emb}$. The mean calibration embedding $\bar{\mathbf{e}}_\text{calib}$ was projected through each layer's gate: $\hat{\mathbf{g}}^{(l)} = \bar{\mathbf{e}}_\text{calib} W_\text{gate}^{(l)}$.

While $\hat{\mathbf{g}}^{(l)}$ provided reasonable alignment in layers 0--2, it diverged completely in deeper layers ($l \ge 10$). In a 30-layer transformer, hidden representations undergo non-linear rotations and multi-head attention projections that shift their geometric trajectory into subspaces orthogonal to $W_\text{emb}$. Static embedding projections misidentified 75\% of active experts in layers 15--29. Pruning based on this proxy removed active deep experts and preserved dormant ones.

\textbf{Remedy:} Mandatory forward-pass profiling via \texttt{llama-imatrix}, ensuring expert selection reflects true hidden-state routing.

\subsection{Failure Mode 2: Cascaded Normalization Suppression}
\textbf{Symptom:} Syntactically formed prose with complete factual emptiness and catastrophic loss of reasoning ability.

\textbf{Root Cause:} When reducing the candidate pool from 128 to 48 experts, the router softmax denominator decreases. To compensate for inflated attention weights, an intermediate implementation applied an attenuation factor of $0.795\times$ to \texttt{ffn\_down\_exps.scale}.

Because this multiplier was applied across 30 consecutive transformer blocks, the attenuation compounded exponentially:
\begin{equation}
(0.795)^{30} \approx 0.00104
\end{equation}
The FFN branch's contribution to the residual stream was effectively suppressed by a factor of 1,000. Each block became almost purely an attention passthrough, stripping the network of its factual retrieval and reasoning capacity.

\textbf{Remedy:} Scales must be preserved at exactly $1.0\times$. At runtime, llama.cpp re-evaluates the softmax across only the surviving expert logits, maintaining proper relative dynamic range automatically.

\subsection{Failure Mode 3: DRY Sampling Interference}
\textbf{Symptom:} Severe code syntax corruption, broken indentation, and missing identifiers in Python outputs.

\textbf{Root Cause:} The interactive benchmarking harness included a Don't-Repeat-Yourself sampler flag: \texttt{--dry-allowed-length 2}, intended to penalize repetitive loops.

However, standard programming syntax relies heavily on 2-character repeating n-grams (\texttt{def}, \texttt{if}, \texttt{in}, \texttt{+=}, \texttt{==}, \texttt{[0}). The DRY sampler severely penalized valid Python keywords, forcing the model to select low-probability corrupted tokens.

\textbf{Remedy:} Removing the DRY sampler and relying on standard nucleus sampling ($\tau=0.4$, $p=0.9$, repetition penalty $1.1$) restored flawless code syntax.

\section{Experimental Results}

\subsection{Experimental Setup}
All benchmarks were performed on a consumer laptop workstation:
\begin{itemize}
\item \textbf{Hardware:} NVIDIA GeForce RTX 4050 Laptop GPU (6141 MiB VRAM), Intel Core i7, 16GB DDR5 system RAM.
\item \textbf{Inference Engine:} llama.cpp (commit \texttt{cbb7d52ec}, CUDA 12).
\item \textbf{Base Model:} Gemma-4-26B-A4B-it-UD-IQ2\_XXS (9.92 GB, 128 experts per layer, 30 layers).
\item \textbf{Compressed Model:} Gemma-4-26B-A4B-Sub5GB-\MAPLE{} (4.75 GB, 48 experts per layer).
\end{itemize}

\subsection{Hardware Throughput and Memory Residency}
Table~\ref{tab:hardware} compares execution telemetry between the original 128-expert base model and the 48-expert \MAPLE{} artifact.

\begin{table}[ht]
\centering
\caption{System Performance on RTX 4050 6GB Laptop GPU}
\label{tab:hardware}
\resizebox{\columnwidth}{!}{%
\begin{tabular}{lcccc}
\toprule
\textbf{Model Container} & \textbf{Disk Size} & \textbf{VRAM Alloc} & \textbf{Layers} & \textbf{Speed} \\
\midrule
Original 128-Exp & 9.92 GB & 6,141 MiB* & 16 / 30 & 6.9 TPS \\
\MAPLE{} 48-Exp & \textbf{4.75 GB} & \textbf{4,180 MiB} & \textbf{30 / 30} & \textbf{45.4 TPS} \\
\midrule
\textbf{Improvement} & \textbf{2.09$\times$ smaller} & \textbf{-32\% VRAM} & \textbf{100\% In-VRAM} & \textbf{6.57$\times$ faster} \\
\bottomrule
\end{tabular}%
}
\end{table}
\footnotesize{*The original model overflowed VRAM, requiring 14 layers to be paged across PCIe into system RAM.}

The 6.57$\times$ speedup represents a direct consequence of eliminating PCIe bus transfers. When 100\% of model tensors reside in GPU VRAM, token generation runs at full memory bandwidth, delivering 45.4 TPS.

\subsection{Effective Parameter Bitrate}
Gemma~4~26B-A4B has $25.6 \times 10^9$ total parameters. With the \MAPLE{} container occupying 4,745,987,168 bytes:
\begin{equation}
\text{Effective BPW} = \frac{4{,}745{,}987{,}168 \times 8}{25.6 \times 10^9} \approx \mathbf{1.48} \text{ bits/parameter}
\end{equation}
This achieves a sub-1.5 BPW effective representation without applying sub-2-bit arithmetic quantization.

\subsection{Task-Level Semantic Fidelity}
We evaluated qualitative output fidelity across four benchmark tasks representing the target deployment domain:
\begin{itemize}
\item \textbf{T1 (Fibonacci Code Synthesis):} Evaluates algorithmic syntax, loop termination, and list accumulation.
\item \textbf{T2 (Binary Search Tree):} Evaluates class structure, insertion logic, and node traversal.
\item \textbf{T3 (Rayleigh Scattering):} Evaluates scientific concept explanation and physical equation citation.
\item \textbf{T4 (Syllogistic Logic):} Evaluates multi-step transitive deduction and formal reasoning.
\end{itemize}

\begin{table}[ht]
\centering
\caption{Task-Level Semantic Keyword Fidelity}
\label{tab:fidelity}
\resizebox{\columnwidth}{!}{%
\begin{tabular}{llccc}
\toprule
\textbf{Task ID} & \textbf{Domain} & \textbf{Keywords} & \textbf{Matched} & \textbf{Fidelity} \\
\midrule
T1: Fibonacci & Code Synthesis & 4 & 4/4 & 100\% \\
T2: BST Search & Data Structures & 4 & 3/4* & 75\% \\
T3: Rayleigh & Physics Deduction & 4 & 4/4 & 100\% \\
T4: Syllogism & Formal Logic & 4 & 4/4 & 100\% \\
\midrule
\textbf{Overall Mean} & & \textbf{16} & \textbf{15/16} & \textbf{93.8\%} \\
\bottomrule
\end{tabular}%
}
\end{table}

In Task T2, the partial 75\% score occurred solely due to generation truncation at 350 tokens, cutting off the \texttt{class Node} definition. Relaxing max tokens to 500 resulted in 100\% completion.

\subsection{Out-of-Domain Generalization and Limitations}
Evaluating perplexity on a general, out-of-domain validation text (\texttt{benchmark\_val.txt}) revealed a perplexity increase from 98.9 (original) to 572.4 (\MAPLE{}). This trade-off is fundamental: \MAPLE{} is explicitly a domain-targeted structural distillation. Users requiring broad multilingual or multi-domain support should retain a higher expert budget ($k' \ge 64$). For edge devices where memory is strictly constrained, \MAPLE{} provides an optimal Pareto trade-off.

\section{Conclusion}
\MAPLE{} demonstrates that fine-grained MoE LLMs contain substantial structural redundancy that can be excised post-hoc. By profiling empirical forward-pass activation manifolds, we compress a 26-billion-parameter MoE model to 1.48 effective bits per parameter, unlocking 45.4 TPS inference on commodity 6~GB laptop GPUs with 93.8\% domain fidelity.

\begin{thebibliography}{00}
\bibitem{jiang2024mixtral} A. Q. Jiang \textit{et al.}, ``Mixtral of Experts,'' \textit{arXiv:2401.04088}, 2024.
\bibitem{dai2024deepseekmoe} D. Dai \textit{et al.}, ``DeepSeekMoE: Towards Ultimate Expert Specialization,'' \textit{arXiv:2401.06066}, 2024.
\bibitem{team2025gemma} Gemma Team, ``Gemma 3 Technical Report,'' \textit{arXiv:2503.19786}, 2025.
\bibitem{frantar2022gptq} E. Frantar \textit{et al.}, ``GPTQ: Accurate Post-Training Quantization,'' \textit{arXiv:2210.17323}, 2022.
\bibitem{lin2024awq} J. Lin \textit{et al.}, ``AWQ: Activation-aware Weight Quantization for On-Device LLM Compression,'' in \textit{Proc. MLSys}, 2024.
\bibitem{dettmers2023spqr} T. Dettmers \textit{et al.}, ``SpQR: A Sparse-Quantized Representation for Near-Lossless LLM Weight Compression,'' \textit{arXiv:2306.03078}, 2023.
\bibitem{frantar2023sparsegpt} E. Frantar and D. Alistarh, ``SparseGPT: Massive Language Models Can Be Accurately Pruned in One Shot,'' in \textit{Proc. ICML}, 2023.
\bibitem{chen2022task} T. Chen \textit{et al.}, ``Task-specific Expert Pruning for Sparse Mixture-of-Experts,'' \textit{arXiv:2206.00277}, 2022.
\bibitem{fedus2022switch} W. Fedus \textit{et al.}, ``Switch Transformers,'' \textit{JMLR}, 2022.
\bibitem{lepikhin2021gshard} D. Lepikhin \textit{et al.}, ``GShard: Scaling Giant Models with Conditional Computation,'' in \textit{Proc. ICLR}, 2021.
\bibitem{llama_cpp} G. Gerganov \textit{et al.}, ``llama.cpp: LLM inference in C/C++,'' \url{https://github.com/ggerganov/llama.cpp}, 2023.
\end{thebibliography}

\end{document}
